% Appendix A: Simulation details
\section{Simulation details}
\label{app:simulation}

\subsection{Configuration}

The simulation is configured via YAML files in \texttt{configs/}:
\begin{itemize}
  \item \texttt{base.yaml}: simulation parameters, kernel choices, cutoffs,
        initial conditions.
  \item \texttt{landscapes.yaml}: well positions $(\mu, V_0)$, basin physics
        vectors $\boldsymbol{\theta}$.
  \item \texttt{kernels.yaml}: kernel names, $\ell$ values, distance types.
\end{itemize}

\subsection{Numerical parameters}

\begin{itemize}
  \item $n_{\rm trajectories} = 10,\!000$ (main run).
  \item Stress test: $100$ runs ($5 \times \sigma \times 20$ seeds),
        $n_{\rm trajectories} = 2,\!000$ per run.
  \item $\sigma_{\rm initial} = 0.5, 1.0, 1.5, 2.0, 3.0$.
  \item Seeds: $1$--$20$.
  \item Observable $\phi_0 = 2.0$.
  \item Primary kernel: Gaussian $\ell = 1.0$.
  \item $e$-folds: $80$, time step $\Delta t = 0.1$.
\end{itemize}

\subsection{Measure computation}

The global census and \ISR measures are computed at each cutoff $N_k$.
KL divergence uses the union of support labels (vacua appearing in either
distribution), with zero probability treated as $10^{-12}$ for
logarithmic stability.
The 1-Wasserstein distance is computed over the cumulative distribution
function of vacuum probabilities.

\subsection{Kernel-weighted sector averaging}

For each $\ell \in \{0.1, 0.25, 0.5, 1.0, 2.0\}$, the aggregate effective
coefficient $c_{\rm eff}(\ell)$ is computed at the final cutoff using
Eq.~\eqref{eq:effective_flow}.
Per-vacuum coefficients are also reported in the output data.
